\documentclass[11pt, letterpaper]{article}
\usepackage{float}
\usepackage{graphicx}
\usepackage{amsmath}
\usepackage[colorlinks=true, allcolors=blue]{hyperref}
\usepackage{fullpage}
\usepackage[margin=1in]{geometry}
\usepackage{booktabs}
\usepackage[labelfont=bf]{caption}
\usepackage{subcaption}
\usepackage{setspace}
\usepackage{parskip}
\usepackage{array}
\usepackage{microtype}
\usepackage{adjustbox}

\setlength{\parindent}{.25in}
\setlength{\parskip}{8pt}

\begin{document}

\title{Rough Draft}
\author{Jood Alaswad, Arnol Garcia, Marissa Miller, Brandon Thacker}
\date{\today}
\maketitle

\tableofcontents
\newpage
\setcounter{page}{1}

\addcontentsline{toc}{section}{Abstract}
\begin{abstract}
Airline cost efficiency is a central issue in the U.S. domestic airline industry, where carriers must balance cost control, operational reliability, and competitive service quality. In recent years, these trade offs have become even more important as airlines have faced major disruptions, shifting demand, and changing business conditions. This study examines how airline cost efficiency relates to operational performance across U.S. domestic carriers, with particular attention to differences across airline business models and across pre-COVID, COVID, and post-COVID periods. Using Bureau of Transportation Statistics T-100, Form 41 Schedule P-1.2, and On-Time Performance data, the project constructs a merged carrier-quarter dataset combining operational, financial, and service-related measures from 2016 through 2024.

The results show that the COVID-19 period functioned as a major structural break in passenger traffic, revenues, expenses, and average net income. Legacy carriers generally exhibited higher route-level utilization than Low-Cost and Regional carriers, while business-model differences were also visible in financial dispersion and service outcomes. At the same time, lower cost per passenger did not map cleanly onto stronger service performance, suggesting that cost minimization alone does not explain reliability outcomes. Delay-cause patterns and clustering analysis further indicate that operational performance varies across carriers in ways that are not fully captured by cost measures alone.

Overall, the findings suggest that airline performance depends not only on leaner cost structures, but also on network design, traffic density, and operational resilience. By integrating traffic, financial, and service datasets into a single carrier-quarter framework, this project provides a broader view of how airline business models perform under both normal and disrupted operating conditions. Future work should extend the analysis with formal regression results and examine how changing market conditions continue to reshape the relationship between cost efficiency and long-run resilience.
\end{abstract}

\section{Introduction and Background}
\subsection{Introduction}
The United States airline industry operates in a highly competitive and cost-sensitive environment, where firms must continuously balance operational efficiency, cost control, and service performance. Airlines face persistent pressure from fluctuating fuel prices, labor costs, and cyclical demand, requiring strategic decisions that optimize both financial and operational outcomes \cite{Barbot2008,Gillen2005,Fan2016}. While cost minimization is often treated as a central objective, its relationship with operational performance remains complex and not fully understood.\\
\indent A key tension in the industry lies in the differing advantages of airline business models. Low-cost carriers are commonly perceived as more efficient due to streamlined operations and lower cost structures, whereas legacy carriers compete through broader networks, connectivity, and higher quality service offerings \cite{Barbot2008,Fan2016,Gillen2005}. This raises an important question: whether cost reduction strategies enhance operational performance or introduce trade offs that negatively affect service quality.\\
\indent Prior research demonstrates that airline efficiency is shaped by multiple operational and strategic factors, including aircraft utilization, load factor, labor productivity, and network design \cite{Caves1984,Banker1993,Zuidberg2014}. At the same time, performance indicators such as delays, cancellations, and reliability introduce additional complexity, as they may reflect both inefficiencies and intentional network design choices \cite{MayerSinai2002,Lee2023}. Furthermore, evolving revenue structures, particularly the growth of ancillary services, have expanded the concept of airline performance beyond traditional measures such as ticket pricing and passenger volume \cite{Zhou2020,Moghadasnian2025}.\\
\indent Given these dynamics, this study examines the relationship between airline cost efficiency and operational performance. It builds on existing literature to evaluate how business models, cost structures, and service metrics interact, while also considering how these relationships evolved across pre-COVID, COVID, and post-COVID periods \cite{FontanetPerez2022,Gudmundsson2021}. By doing so, the analysis aims to determine whether cost efficiency and service quality function as complementary objectives or reflect inherent trade offs within the airline industry.

\subsection{Background}
The modern airline industry has been fundamentally shaped by deregulation, liberalization, and increasing competition, particularly following the United States Airline Deregulation Act of 1978. Deregulation intensified competitive pressures and forced airlines to prioritize cost efficiency while maintaining service quality and network coverage \cite{Borenstein1992, Gillen2005}. This shift also enabled the rapid emergence of low-cost carriers, which introduced alternative operating strategies that challenged traditional airline models.\\
\indent Airline business models represent a central component of industry structure. Low-cost carriers typically operate low complexity networks with high aircraft utilization, simplified fleets, dense seating configurations, and limited service offerings. In contrast, legacy carriers rely on hub and spoke systems with broader connectivity, and differentiated services \cite{Barbot2008, Gillen2005, OConnell2005}. These structural differences influence how airlines allocate resources, manage capacity, and pursue efficiency, making business model distinctions critical to understanding cost performance.\\
\indent Prior research finds that cost efficiency is closely tied to operational factors such as traffic density, aircraft utilization, and network design. In particular, economies of density, rather than economies of scale alone, play a significant role in reducing airline costs \cite{Caves1984, Zuidberg2014}. Historical empirical evidence suggests that low-cost carriers are generally more efficient than full-service carriers, although this relationship is influenced by geography, labor productivity, and firm size \cite{Barbot2008}. More recent studies emphasize the importance of accounting for heterogeneity across airlines, cautioning against treating carriers as homogeneous entities in efficiency analysis \cite{Li2025}.\\
\indent Operational performance adds another layer of complexity. Airlines regularly experience disruptions such as delays and cancellations, which affect reliability and customer satisfaction. However, not all delays are indicative of inefficiency. Congestion at hub airports can emerge as a rational outcome of network optimization, where airlines cluster flights to maximize connectivity \cite{MayerSinai2002}. This suggests that trade offs between efficiency and service performance may be inherent to airline operations rather than solely the result of poor management. Additional research highlights the importance of disruption mitigation and operational resilience, emphasizing the balance between minimizing costs and maintaining service reliability \cite{Lee2023}.\\
\indent Airline cost and revenue structures have also evolved significantly over time. Increasing reliance on ancillary revenue streams, such as baggage fees and seat selection, has expanded how airline performance is evaluated \cite{Zhou2020}. As a result, understanding airline efficiency requires a broader perspective that incorporates cost structures, revenue strategies, and operational outcomes.

\section{Methods}
\subsection{Methods: Data Sources, Preparation, and Cleaning}
This project is based on three datasets from the Bureau of Transportation Statistics (BTS), enabling cross-analysis of airline traffic, financial outcomes, and service performance.\\
\indent The \textbf{T-100 Domestic Segment} dataset \cite{bts2026t100} provides route-level traffic information, such as the amount of passengers, freight, and mail enplaned as well as the distance between origin and destination airports. A route key for each observation was created using origin and destination airport identifiers for later determining how many unique routes a carrier performs. Numerical data were coerced into numeric for smooth aggregation and creating calculated variables later on. Route-level observations were aggregated after assigning each with a composite match key formed from a carrier’s unique code, year, and quarter.\\
\indent The \textbf{Form 41 Schedule P-1.2} financial dataset \cite{bts2026f41} contains quarterly profit and loss statements for large U.S. carriers. It includes net income and breaks down operating revenues and operating expenses, making it relevant for our focus on operating performance. The P-1.2 financial data were processed and aggregated similarly to the T-100 data.\\
\indent The \textbf{On-Time Performance (OTP)} dataset \cite{bts2026otp}, sourced from Kaggle \cite{kaggle_bts2026otp} but derived from BTS records, details monthly arrival delay information, broken down by airline and airport. It originally contains 398,233 observations from 2003 through 2024, but for this analysis all pre-2016 data is filtered out. Certain variables in the BTS data were not included in the Kaggle version, most notably it did not contain unique carrier identifiers. General carrier codes can be re-used by unrelated carriers over time and general carrier names can have different naming conventions, causing carrier mismatches if merged as-is with the other data. To ensure consistency and accuracy throughout the analysis, carrier identifiers were standardized according to the unique carrier names and codes established by BTS. The OTP data were then aggregated similarly to the T-100 and P-1.2 data.\\
\indent Once each dataset was cleaned and individually aggregated, they were merged into one dataset using a left join with T-100 as the base to preserve all operational records. Observations were categorized as \textit{Pre-COVID} (2016-2019), \textit{COVID} (2020-2021), or \textit{Post-COVID} (2022-2024). Carriers were classified into four main groups: \textit{Legacy}, \textit{Low-Cost}, \textit{Cargo}, and \textit{Regional}. The group assigned to each carrier was based on industry classification definitions; carriers that did not clearly align with any one type were assigned to '\textit{other}'. Cargo-only operators are largely absent from the OTP data, so service-quality findings should be interpreted with that limitation in mind. Finally, multiple variables were engineered to create easily interpretable metrics. These include operating margin, net margin, economics per passenger and route, arrival performance rates, and delay cause ratios, among several other measures.

\begin{figure}[H]
    \centering
    \includegraphics[width=\textwidth]{data_shape_column_names.png}
    \caption{Data shape of datasets before and after cleaning/aggregation, along with column names of the final merged dataset. Upper-case variables represent original data, and lower-case variables represent engineered features.}
    \label{fig:shape}
\end{figure}

\subsection{Methods: Exploratory Data Analysis and Visualization}
After cleaning and merging the datasets, exploratory data analysis was used to understand the structure of the carrier-quarter dataset and identify the most relevant patterns for later modeling. Descriptive statistics were calculated for key financial, operational, and service variables, with particular attention to the distribution of passengers, freight, route count, operating revenue, operating expenses, net income, operating margin, cost per passenger, on-time rate, and cancellation rate.

A range of visualizations was used to examine structure and relationships in the data. Yearly trend plots were used to evaluate changes in passenger traffic, freight, revenues, expenses, net income, and on-time rate over time. Because the COVID-19 pandemic represented a major disruption to airline operations, the period was divided into pre-COVID (2016--2019), COVID (2020--2021), and post-COVID (2022--2024) for comparison. Boxplots and violin plots were used to compare cost, margin, and service-performance measures across airline business models, including Legacy, Low-Cost, Cargo, and Regional carriers. Scatterplots were used to examine whether lower cost per passenger was associated with stronger service outcomes, particularly on-time rate and cancellation rate. A correlation heatmap was also constructed to summarize pairwise relationships among operational, financial, and service variables.

These exploratory steps were not used as final evidence on their own. Instead, they were used to identify broad patterns, highlight differences across carrier types and periods, and motivate the more formal statistical and machine learning methods described below.

\subsection{Methods: Statistical and Inferential Analysis}
To address the project’s inferential component, ordinary least squares regression was selected as the primary statistical framework. This was preferred over isolated group mean tests because the research questions involve multiple predictors, multiple outcomes, and important subgroup differences across carrier type and COVID period. The main goal of the regression analysis is therefore to estimate whether cost structure, utilization, and business-model indicators are associated with operational and service outcomes while accounting for time period and network structure.

For the first research question, the project uses average passengers per route as a proxy for route-level utilization, since the version of T-100 used here does not include seats or departures performed and therefore does not permit direct calculation of load factor. A baseline regression specification is:

\begin{align*}
\text{PassengersPerRoute}_{it} =\;& \beta_0
+ \beta_1 \text{CostPerPassenger}_{it}
+ \beta_2 \text{AvgDistancePerRoute}_{it} \\
&+ \beta_3 \text{RouteCount}_{it}
+ \beta_4 \text{CarrierType}_{i}
+ \beta_5 \text{Period}_{t}
+ \epsilon_{it}
\end{align*}

For the second research question, service outcomes are modeled using on-time rate and cancellation rate as dependent variables. Baseline specifications include:

\begin{align*}
\text{OnTimeRate}_{it} =\;& \beta_0
+ \beta_1 \text{CostPerPassenger}_{it}
+ \beta_2 \text{CarrierType}_{i} \\
&+ \beta_3 \text{Period}_{t}
+ \beta_4 \text{AvgDistancePerRoute}_{it}
+ \epsilon_{it}
\end{align*}

\begin{align*}
\text{CancellationRate}_{it} =\;& \beta_0
+ \beta_1 \text{CostPerPassenger}_{it}
+ \beta_2 \text{CarrierType}_{i} \\
&+ \beta_3 \text{Period}_{t}
+ \beta_4 \text{AvgDistancePerRoute}_{it}
+ \epsilon_{it}
\end{align*}

These models are intended to estimate associations rather than establish causal effects. Coefficient direction and magnitude are interpreted alongside the exploratory findings.

\subsection{Methods: Machine Learning and Unsupervised Analysis}
In addition to regression-based analysis, the project uses Random Forest and K-means clustering as supplementary methods. A Random Forest model was applied to identify which operational, financial, and service variables were most important in predicting quarterly net income. This approach is useful for ranking feature importance and for identifying nonlinear patterns that may not be captured well by simple linear models. However, because some financial variables are closely related to net income by construction, the Random Forest results are interpreted cautiously and are used primarily for exploratory prioritization rather than causal inference.

K-means clustering was used to identify recurring airline operating profiles based on financial, operational, and service variables. Before clustering, variables were standardized so that large-scale revenue or passenger values would not dominate the distance metric. The elbow method was used to guide the selection of the number of clusters. The resulting cluster profiles were compared against the pre-defined carrier-type categories to determine whether airlines naturally separate into distinct operating groups such as large passenger carriers, mixed mid-scale carriers, distressed low-passenger groups, or cargo-heavy carriers.

Together, these methods allow the project to combine descriptive analysis, inferential statistics, and machine learning in a way that aligns with the project’s research questions and the broader goals of the capstone.

\section{Results}

\subsection{Merged Sample Coverage and Descriptive Results}
The final merged airline dataset was constructed at the carrier-quarter level using T-100 Domestic Segment, Form 41 Schedule P-1.2, and On-Time Performance data. The T-100 base contained 4,440 carrier-quarter observations. Of these, 1,825 matched Form 41 financial observations, 739 matched OTP service observations, and 681 matched all three datasets simultaneously. These coverage differences are important because they determine which analyses can be conducted on the full operational sample and which must rely on smaller matched subsets.

\begin{table}[H]
\centering
\caption{Merged Sample Coverage}
\label{tab:coverage}
\begin{tabular}{lr}
\toprule
Sample Component & Carrier-Quarter Rows \\
\midrule
T-100 base sample & 4,440 \\
Matched with Form 41 & 1,825 \\
Matched with OTP & 739 \\
Matched across all three datasets & 681 \\
\bottomrule
\end{tabular}
\end{table}

Descriptive statistics show that the merged dataset is highly heterogeneous. Passenger volume, freight, distance, and route count vary substantially across observations, and financial measures such as operating revenue, operating expense, net income, cost per passenger, and revenue per passenger are strongly right-skewed. This confirms that the airline industry cannot be treated as homogeneous and that cost-performance relationships must be interpreted in the context of business model and network structure.

\subsection{COVID-19 as a Structural Break}
One of the clearest findings in the data is that the COVID-19 period represents a major structural break in U.S. domestic airline operations. Passenger volume increased through 2019, collapsed sharply in 2020, and then recovered over the following years, though not uniformly across carrier types. Freight behaved differently, remaining more stable and even increasing during parts of the COVID period. Operating revenues and operating expenses both shifted sharply around 2020, while average quarterly net income became deeply negative.

On-time rate improved temporarily during the COVID period. However, this should not be interpreted as a simple improvement in operational efficiency. A more plausible explanation is that reduced schedule density lowered congestion and delay propagation. Taken together, these results support the decision to analyze pre-COVID, COVID, and post-COVID periods separately rather than pooling all years together.

\begin{figure}[H]
    \centering
    \includegraphics[width=\textwidth]{fig0_yearly_trends.png}
    \caption{Yearly industry trends in passenger volume, freight, operating revenue, operating expense, average quarterly net income, and on-time rate. The figure shows a clear structural disruption during the COVID period.}
    \label{fig:yearly_trends}
\end{figure}

\subsection{Carrier-Type and Utilization Differences}
Carrier type differences were also pronounced. Using average passengers per route as a utilization proxy, Legacy carriers consistently showed higher route-level traffic concentration than Low-Cost and Regional carriers, particularly before COVID. All groups experienced declines during the pandemic, but their recovery patterns differed. This suggests that network structure and traffic density matter alongside cost when evaluating operational efficiency.

\begin{figure}[H]
    \centering
    \includegraphics[width=\textwidth]{fig_loadfactor.png}
    \caption{Average passengers per route by carrier type and period. Legacy carriers consistently operate with higher route-level traffic concentration than Low-Cost and Regional carriers.}
    \label{fig:loadfactor}
\end{figure}

Comparisons of financial and operational KPIs also revealed differences across business models. Legacy carriers showed wider distributions in operating margin and on-time rate, while Low-Cost carriers were more tightly clustered. Cargo carriers exhibited extreme passenger-based ratios such as cost per passenger and revenue per passenger, but these values are driven by near-zero passenger denominators and therefore should not be interpreted in the same way as passenger-carrier metrics.

\begin{figure}[H]
    \centering
    \includegraphics[width=\textwidth]{fig2_carrier_kpis.png}
    \caption{Financial and operational KPI distributions by carrier type. Cargo carriers exhibit extreme passenger-based ratios because passenger denominators are very small, while Legacy and Low-Cost carriers show clearer differences in spread and central tendency.}
    \label{fig:carrier_kpis}
\end{figure}

\subsection{Cost and Service Performance}
The project’s central question concerns whether lower-cost structures align with better service performance. The exploratory results suggest that the relationship is not simple. Scatterplots comparing cost per passenger with on-time rate and cancellation rate did not show a strong linear relationship. Carriers with similar cost levels still exhibited different service outcomes, indicating that business model and network structure likely shape service performance more strongly than cost level alone.

Delay-cause analysis provides additional context. Late-aircraft delay made up the largest share of delay minutes across all carrier types, while NAS delay was relatively larger for Legacy and Regional carriers and carrier delay was proportionally larger for Low-Cost carriers. This suggests that a weaker on-time rate can reflect different underlying operational mechanisms depending on the airline type.

\begin{figure}[H]
    \centering
    \includegraphics[width=0.85\textwidth]{fig4_delay_causes.png}
    \caption{Delay-cause composition by carrier type. Late-aircraft delay is the largest component across groups, while NAS and carrier delay vary by airline type.}
    \label{fig:delay_causes}
\end{figure}

\subsection{Correlation Structure and Cluster Profiles}
The correlation analysis showed that financial performance is not strongly explained by any single operational variable on its own. Operational scale variables were only weakly correlated with net income, which suggests that profitability depends on combinations of cost structure, network design, and service performance rather than on traffic volume alone.

\begin{figure}[H]
    \centering
    \includegraphics[width=\textwidth]{fig7_correlation.png}
    \caption{Correlation heatmap across operational, financial, and service metrics. Strong relationships exist among some financial variables, but no single operational measure fully explains financial performance.}
    \label{fig:correlation}
\end{figure}

K-means clustering produced four distinct operating profiles. One cluster represented large-scale passenger carriers with high traffic volume and relatively stronger overall performance. A second represented mixed mid-scale carriers with moderate performance. A third represented extreme distressed or outlier observations with very negative margins and almost no passenger base. A fourth appeared cargo-dominant, with extremely high passenger-based ratios driven by very low passenger denominators. These results reinforce the idea that simple carrier-type labels are useful but not sufficient to capture all of the variation in the data.

\begin{figure}[H]
    \centering
    \includegraphics[width=\textwidth]{fig_cluster_profiles.png}
    \caption{K-means cluster profiles based on operational, financial, and service variables. The four-cluster solution identifies large passenger carriers, mixed mid-scale carriers, distressed/outlier observations, and cargo-dominant profiles.}
    \label{fig:clusters}
\end{figure}

\section{Discussion}

The results point to several broader themes. First, airline cost efficiency cannot be reduced to simple size measures. The utilization analysis suggests that route-level traffic concentration differs substantially across business models, which is consistent with prior work emphasizing density economies, traffic concentration, and route structure as major determinants of airline cost performance \cite{Banker1993, Caves1984, Li2025, Zuidberg2014}.

Second, the relationship between cost and service quality is not straightforward. The project’s EDA and early results suggest that lower cost per passenger does not reliably produce better service outcomes. On-time performance and cancellation patterns appear to be shaped by congestion, delay propagation, network structure, and carrier-specific operating models rather than by cost minimization alone \cite{Lee2023, MayerSinai2002}. This is important because it challenges any simple assumption that lower-cost carriers are always operationally superior.

Third, the results support the idea that business-model differences remain central to understanding airline performance. Legacy, Low-Cost, Regional, and Cargo carriers differ not just in cost levels, but also in utilization, route design, financial dispersion, and service-performance patterns. This is broadly consistent with prior literature on airline heterogeneity and revenue or efficiency differences across business models \cite{Barbot2008, Gillen2005, Nguyen2022, Li2025}.

This project also contributes to the literature by merging three types of airline information, operational, financial, and service-related, into a single carrier-quarter framework. Even though overlap across datasets is incomplete, that integrated perspective makes it possible to examine airline performance more holistically than studies focused on only one outcome type.

The project also raises ethical considerations. A cost-efficiency framework can be misused if it is interpreted as recommending aggressive cost minimization without regard to service reliability, labor strain, or regional access. What appears efficient at an aggregate level may increase fragility or worsen service for certain categories of travelers or smaller markets. There is also an ethical risk in treating exploratory or predictive results as objective rankings of airlines when the data are incomplete and the relationships are not causal.

As a contemporary development outside the study period, Spirit Airlines announced on May 2, 2026 that it had begun an orderly wind-down of operations effective immediately, cancelling all flights and ending customer service \cite{reuters_spirit_shutdown2026,ap_spirit_shutdown2026,spirit_shutdown2026}. Although this event falls outside the 2016--2024 analytical window used in this study, it provides useful context for the project’s broader findings. Spirit’s collapse suggests that while low-cost carriers may operate with leaner structures, those same models may also remain vulnerable to financial stress, external shocks, and thin margins. In that sense, the shutdown reinforces the project’s broader conclusion that lower cost does not automatically imply stronger long-run resilience or performance.

\section{Conclusion}

This project examined the relationship between airline cost efficiency and operational performance in the U.S. domestic airline industry by combining T-100 traffic data, Form 41 financial data, and On-Time Performance data at the carrier-quarter level. Across the analysis, the most consistent result is that cost, efficiency, and service quality do not move together in a simple way.

Legacy carriers generally operate with higher route-level traffic concentration, while Low-Cost, Regional, and Cargo carriers show distinct patterns in utilization, financial performance, and service quality. The COVID-19 period represented a clear structural break, affecting both financial and operational outcomes across the industry. More broadly, the results suggest that airline performance depends not only on cost minimization, but also on business model, network structure, and operational resilience.

The bottom line is that lower-cost strategies may produce some efficiency advantages, but they do not automatically guarantee better service or stronger long-term performance. Recent industry developments, including Spirit Airlines’ May 2026 shutdown, further suggest that cost minimization alone is not enough to ensure resilience or continued market viability \cite{reuters_spirit_shutdown2026,spirit_shutdown2026}.

\clearpage
\addcontentsline{toc}{section}{References}
\bibliographystyle{abbrv}
\bibliography{RoughDraftRefs}

\clearpage
\appendix
\section{Project Code}
GitHub Repository: \url{https://github.com/mnmill21/DAT490_Project}

\end{document}